\chapter{Competências Digitais e Informática Aplicada ao Setor Público}

\section{Cultura digital, competência digital e cidadania}

Cultura digital é o conjunto de práticas, valores, linguagens e formas de interação que surgem da presença de tecnologias digitais na vida social, profissional e institucional. Competência digital é mais específica: corresponde à capacidade de utilizar tecnologias de forma crítica, segura, ética e eficiente para localizar informação, comunicar, produzir conteúdo, resolver problemas e exercer direitos.

No setor público, competência digital não se limita ao domínio operacional de programas. Ela envolve compreender como dados, sistemas, redes e plataformas modificam processos administrativos, prestação de serviços, transparência, participação social e proteção de direitos.

\begin{teoria}[quatro dimensões da competência digital]
Uma forma útil de organizar o tema é separar quatro dimensões:
\begin{enumerate}
\item \textbf{informacional}: localizar, avaliar, organizar e interpretar informação;
\item \textbf{tecnológica}: operar ferramentas, redes, aplicações e dispositivos;
\item \textbf{comunicacional}: produzir e compartilhar conteúdo adequadamente;
\item \textbf{ética e jurídica}: respeitar privacidade, segurança, direitos autorais, transparência e responsabilidades.
\end{enumerate}
\end{teoria}

A Política Nacional de Educação Digital, instituída pela Lei nº 14.533/2023, estrutura-se em eixos como inclusão digital, educação digital escolar, capacitação e especialização digital e pesquisa e desenvolvimento em TIC. A lei também relaciona educação digital a letramento digital, lógica, algoritmos, programação, ética, letramento midiático e cidadania digital.

\section{Pensamento computacional}

Pensamento computacional é uma forma de estruturar problemas para que possam ser compreendidos e resolvidos de maneira sistemática, com ou sem uso direto de computador.

Quatro ideias aparecem com frequência:
\begin{itemize}
\item \termo{decomposição}: dividir problema complexo em partes menores;
\item \termo{reconhecimento de padrões}: identificar regularidades entre casos;
\item \termo{abstração}: selecionar características relevantes e ignorar detalhes irrelevantes;
\item \termo{algoritmo}: definir sequência finita e não ambígua de passos para atingir um objetivo.
\end{itemize}

\begin{exemplo}[processo administrativo]
Para automatizar a triagem de processos, pode-se decompor o fluxo em recebimento, validação documental, classificação, encaminhamento e registro. Em seguida, identificam-se regras repetitivas, abstraem-se campos essenciais e descreve-se um algoritmo para decidir o próximo estado de cada processo.
\end{exemplo}

Pensamento computacional não equivale a programação. Programar é uma possível implementação de uma solução; o raciocínio algorítmico pode ser aplicado antes da escolha da linguagem ou mesmo em processos executados manualmente.

\section{Informação, dado e conhecimento}

\termo{Dado} é uma representação elementar de fato, evento ou atributo. \termo{Informação} resulta da organização e interpretação de dados em contexto. \termo{Conhecimento} envolve compreensão que permite explicar, decidir ou agir.

Uma lista de valores de despesas é dado. O total por unidade e por mês já constitui informação. A identificação de um padrão anômalo e sua interpretação à luz do processo de contratação aproxima-se de conhecimento.

Essa distinção é importante porque sistemas podem armazenar grandes volumes de dados sem produzir informação de qualidade. Qualidade depende de atributos como exatidão, completude, consistência, atualidade, unicidade e adequação ao uso.

\section{Internet, Web e endereçamento}

Internet é a infraestrutura global baseada em interconexão de redes e protocolos da família TCP/IP. Web é um serviço executado sobre essa infraestrutura, normalmente utilizando HTTP ou HTTPS.

\subsection{Estrutura de uma URL}

Considere:
\begin{center}
\texttt{https://portal.exemplo.gov.br/processos?id=42\#historico}
\end{center}

Podemos separar:
\begin{itemize}
\item \texttt{https}: esquema/protocolo;
\item \texttt{portal.exemplo.gov.br}: host;
\item \texttt{/processos}: caminho;
\item \texttt{id=42}: parâmetro de consulta;
\item \texttt{historico}: fragmento.
\end{itemize}

O DNS traduz nomes de domínio em informações de endereçamento. O navegador não ``procura a página no DNS''; ele utiliza o DNS para resolver o nome do host e depois estabelece comunicação com o serviço correspondente.

\subsection{HTTP e HTTPS}

HTTP é protocolo de aplicação usado para transferência de recursos na Web. HTTPS corresponde a HTTP protegido por TLS. O TLS fornece propriedades como confidencialidade e integridade do canal e autenticação do servidor por certificados, conforme a configuração.

\begin{atencao}
HTTPS não garante que o conteúdo do site seja verdadeiro, que a organização seja confiável ou que o usuário não esteja diante de uma página maliciosa. Ele protege a comunicação com o servidor identificado pelo certificado; não valida a intenção do conteúdo.
\end{atencao}

\section{Navegadores, cache, cookies e armazenamento local}

Navegadores interpretam HTML, CSS e JavaScript, gerenciam sessões, certificados, cache, cookies, extensões e armazenamento local.

\termo{Cache} mantém cópias de recursos para reduzir tráfego e tempo de carregamento. \termo{Cookie} é pequeno dado associado a um domínio e enviado conforme regras da aplicação; pode manter sessão, preferências ou identificadores. \termo{Local storage} permite persistir dados no navegador sem o mesmo mecanismo de envio automático dos cookies.

Cookies não são, por definição, vírus. Contudo, identificadores podem ser usados para rastreamento e perfilamento. O impacto depende do conteúdo, finalidade e contexto de uso.

O modo privado ou anônimo reduz a persistência local de histórico, cookies e outros dados ao encerrar a sessão, mas não torna a conexão anônima perante provedor, rede institucional, servidor acessado ou mecanismos externos de correlação.

\section{Motores de busca e recuperação de informação}

Motores de busca realizam rastreamento, indexação e ranqueamento. O resultado exibido primeiro não é necessariamente a fonte mais correta; é o resultado que o mecanismo considerou mais relevante segundo seus critérios.

\subsection{Operadores de busca}

Operadores úteis incluem:
\begin{itemize}
\item aspas para expressão exata: \texttt{"tribunal de contas"};
\item \texttt{site:dominio} para restringir domínio;
\item \texttt{filetype:pdf} para buscar determinado formato, quando suportado;
\item sinal de menos para excluir termo em mecanismos que adotem essa sintaxe.
\end{itemize}

\subsection{Fonte primária e fonte secundária}

Fonte primária apresenta o documento, dado ou ato em sua origem: texto legal oficial, acórdão, base estatística, edital, relatório institucional. Fonte secundária interpreta a fonte primária: notícia, comentário, resumo ou artigo.

Em temas normativos, o uso de fonte primária reduz o risco de estudar dispositivo revogado ou interpretação desatualizada.

\section{Letramento midiático e desinformação}

Letramento midiático é a capacidade de acessar, compreender, avaliar e produzir mensagens em diferentes meios. Em ambiente digital, isso exige distinguir evidência, opinião, publicidade, sátira, conteúdo manipulado e informação fora de contexto.

A desinformação pode explorar:
\begin{itemize}
\item título que não corresponde ao conteúdo;
\item imagem verdadeira associada a evento diferente;
\item dado correto apresentado sem denominador ou contexto;
\item gráfico com eixo truncado;
\item falsa autoridade;
\item conteúdo sintético ou editado;
\item rede coordenada de repetição para criar aparência de consenso.
\end{itemize}

\begin{exemplo}[gráfico enganoso]
Se uma taxa sobe de 50 para 52, um gráfico cujo eixo vertical comece em 49 pode representar visualmente um salto enorme. O dado não foi alterado, mas a escala pode induzir interpretação desproporcional. Letramento digital inclui examinar origem, unidade, denominador, período e escala.
\end{exemplo}

\section{Correio eletrônico}

E-mail combina agentes de usuário, servidores e protocolos. SMTP é associado ao envio e transferência de mensagens entre servidores. IMAP permite acesso mantendo mensagens organizadas no servidor. POP3, em configuração clássica, é orientado ao download das mensagens, embora clientes possam manter cópias no servidor.

Os campos mais comuns são:
\begin{itemize}
\item \textbf{Para}: destinatários principais;
\item \textbf{Cc}: cópia visível;
\item \textbf{Cco}: cópia cujo destinatário não é revelado aos demais recipientes na mensagem entregue;
\item \textbf{Assunto}: metadado descritivo;
\item \textbf{Anexo}: arquivo associado à mensagem.
\end{itemize}

Cabeçalhos de e-mail podem conter registros técnicos de encaminhamento. O nome exibido do remetente não prova sua identidade; por isso, golpes exploram domínio semelhante, spoofing e contas comprometidas.

\section{Ferramentas de colaboração e computação em nuvem}

Editores colaborativos, armazenamento em nuvem, calendários, chats e plataformas de reunião permitem coautoria e sincronização entre dispositivos.

\subsection{Permissões}

Três níveis aparecem com frequência:
\begin{itemize}
\item leitura;
\item comentário;
\item edição.
\end{itemize}

Compartilhar ``com qualquer pessoa que possua o link'' é conceitualmente diferente de compartilhar com identidades previamente definidas. O primeiro modelo torna o segredo do link parte relevante do controle de acesso; o segundo depende de autenticação da identidade autorizada.

\subsection{Versionamento e sincronização}

Histórico de versões registra estados anteriores do arquivo e pode permitir restauração. Sincronização mantém cópias coerentes entre dispositivo e serviço remoto, mas não é sinônimo de backup. Uma exclusão ou criptografia maliciosa pode ser sincronizada para outros dispositivos.

\begin{atencao}[sincronização não é backup]
Backup pressupõe capacidade de recuperação a partir de cópia preservada. Serviço de sincronização pode ajudar, especialmente quando possui versionamento, mas sua finalidade principal é manter estados sincronizados, não necessariamente garantir retenção independente contra perda ou ataque.
\end{atencao}

\section{Backup e recuperação}

Backup é cópia destinada à recuperação após exclusão, corrupção, falha, desastre ou incidente de segurança.

\subsection{Tipos de backup}

\begin{table}[ht]
\centering
\small
\begin{tabularx}{\textwidth}{l X X}
\toprule
Tipo & O que copia & Recuperação \\
\midrule
Completo & todo o conjunto selecionado & simples, mas com maior volume de cópia. \\
Incremental & alterações desde o último backup de qualquer tipo & exige cadeia de incrementais após o último completo. \\
Diferencial & alterações desde o último completo & exige o completo e o diferencial mais recente. \\
\bottomrule
\end{tabularx}
\end{table}

Uma prática conhecida como regra 3-2-1 recomenda múltiplas cópias em meios distintos, com pelo menos uma em localização ou domínio de falha separado. Mais importante que memorizar a regra é compreender o princípio: reduzir a chance de um mesmo incidente destruir dados de produção e todas as cópias.

Backup sem teste de restauração produz evidência fraca de recuperabilidade.

\section{Malware e ameaças digitais}

\termo{Malware} é categoria geral de software malicioso.

\begin{table}[ht]
\centering
\small
\begin{tabularx}{\textwidth}{l X}
\toprule
Tipo & Característica principal \\
\midrule
Vírus & associa-se a arquivo ou objeto hospedeiro e executa código malicioso quando ativado. \\
Worm & propaga-se automaticamente entre sistemas explorando rede ou vulnerabilidades. \\
Trojan & apresenta-se como software legítimo ou desejável para induzir execução. \\
Ransomware & bloqueia ou cifra recursos para extorsão, frequentemente combinado com exfiltração. \\
Spyware & coleta informações de forma indevida. \\
Rootkit & busca ocultar presença ou manter acesso privilegiado. \\
\bottomrule
\end{tabularx}
\end{table}

\subsection{Phishing, spear phishing e pharming}

Phishing utiliza engenharia social para induzir a vítima a revelar segredo, executar arquivo ou acessar página maliciosa. Spear phishing é direcionado a vítima ou organização específica. Pharming busca redirecionar o usuário a destino fraudulento, por exemplo mediante manipulação de resolução de nomes ou ambiente local.

Antivírus, firewall e anti-spyware são controles distintos. Nenhum fornece proteção absoluta; segurança depende de camadas complementares como atualização, autenticação forte, filtragem, backup, segmentação e conscientização.

\section{Autenticação e fatores}

Autenticação verifica identidade. Os fatores clássicos incluem:
\begin{itemize}
\item algo que o usuário \textbf{sabe}: senha/PIN;
\item algo que o usuário \textbf{possui}: token, dispositivo ou chave;
\item algo que o usuário \textbf{é}: característica biométrica.
\end{itemize}

MFA exige fatores de categorias distintas. Duas senhas não constituem autenticação multifator, pois ambas pertencem à categoria ``algo que sabe''.

\section{LGPD: fundamentos e conceitos}

A Lei nº 13.709/2018 regula o tratamento de dados pessoais, inclusive em meios digitais, por agentes públicos e privados nas hipóteses de seu âmbito de aplicação. Seu objetivo inclui proteger liberdade, privacidade e livre desenvolvimento da personalidade.

\subsection{Dado pessoal, dado sensível e anonimização}

\termo{Dado pessoal} é informação relacionada a pessoa natural identificada ou identificável. \termo{Dado pessoal sensível} recebe proteção especial e inclui categorias como origem racial ou étnica, convicção religiosa, opinião política, saúde, vida sexual e dados genéticos ou biométricos vinculados a pessoa natural.

Dado anonimizado é aquele que não permite identificar o titular considerando meios técnicos razoáveis e disponíveis, nos termos da lei. Pseudonimização reduz associação direta, mas não equivale necessariamente a anonimização irreversível.

\subsection{Tratamento}

Tratamento é conceito amplo: inclui coleta, produção, recepção, classificação, utilização, acesso, reprodução, transmissão, distribuição, processamento, arquivamento, armazenamento, eliminação e outras operações previstas legalmente.

\subsection{Princípios}

A LGPD estabelece princípios como finalidade, adequação, necessidade, livre acesso, qualidade, transparência, segurança, prevenção, não discriminação e responsabilização/prestação de contas.

\begin{exemplo}[necessidade]
Se um formulário de inscrição precisa apenas confirmar que o candidato é maior de idade, coletar cópia integral de documento com informações não necessárias pode violar o princípio da necessidade, dependendo do contexto e da base jurídica aplicável.
\end{exemplo}

\subsection{Bases legais}

Consentimento é uma das bases legais, não a única. A lei prevê outras hipóteses, como cumprimento de obrigação legal ou regulatória, execução de políticas públicas, exercício regular de direitos, proteção da vida, tutela da saúde, legítimo interesse nas condições legais e proteção do crédito, além de bases específicas para dados sensíveis.

No poder público, o tratamento deve ser examinado à luz da competência e finalidade pública, das regras específicas da LGPD e da legislação que fundamenta a atividade.

\subsection{Controlador, operador e encarregado}

\termo{Controlador} toma as decisões referentes ao tratamento. \termo{Operador} realiza tratamento em nome do controlador. \termo{Encarregado} desempenha as funções previstas na legislação e regulação, servindo, entre outras atribuições, como canal de comunicação.

A qualificação depende da atuação concreta. Um contrato chamar determinada empresa de ``operadora'' não encerra a análise se, na prática, ela decidir autonomamente finalidades e meios relevantes do tratamento.

\section{LAI e transparência}

A Lei nº 12.527/2011 regulamenta o direito de acesso à informação pública. Transparência \termo{ativa} é a divulgação de informação independentemente de solicitação. Transparência \termo{passiva} ocorre quando a Administração responde a pedido de acesso.

A publicidade é regra, mas não é absoluta. Há informações classificadas, sigilos legalmente protegidos e dados pessoais sujeitos a regime próprio. A aplicação conjunta de LAI e LGPD exige distinguir informação sobre a atuação estatal de dados pessoais cuja exposição não seja necessária ou proporcional.

\section{Marco Civil da Internet}

A Lei nº 12.965/2014 estabelece princípios, garantias, direitos e deveres para o uso da Internet no Brasil.

Entre os temas centrais estão liberdade de expressão, privacidade, proteção de dados, neutralidade de rede, registros e responsabilidade dos agentes segundo o regime legal.

\subsection{Neutralidade de rede}

Neutralidade não significa que todos os pacotes terão exatamente a mesma latência. O princípio trata do tratamento isonômico do tráfego nos termos da lei e da regulamentação, admitindo hipóteses legalmente justificadas de discriminação ou degradação.

\section{Interoperabilidade no setor público}

Interoperabilidade é a capacidade de sistemas, organizações ou componentes trocarem informação e utilizarem adequadamente o que foi trocado.

É útil separar dimensões:
\begin{itemize}
\item \textbf{técnica}: protocolos, formatos, APIs e conectividade;
\item \textbf{semântica}: significado compartilhado dos dados;
\item \textbf{organizacional}: processos, responsabilidades e acordos entre instituições;
\item \textbf{jurídica}: autorização e limites para compartilhamento e uso.
\end{itemize}

Dois sistemas podem estar tecnicamente conectados e, ainda assim, não serem semanticamente interoperáveis se o campo ``situação'' possuir códigos e significados incompatíveis.

\section{Processo eletrônico e SEI}

Processo eletrônico substitui ou complementa fluxos baseados em papel por documentos, assinaturas, tramitação, metadados e registros digitais. Seus benefícios incluem rastreabilidade, busca, acesso simultâneo e redução de movimentação física, mas dependem de governança documental, controle de acesso e preservação.

O Sistema Eletrônico de Informações (SEI) é amplamente utilizado na Administração Pública para criação e tramitação de processos e documentos eletrônicos. Conceitualmente, o candidato deve compreender unidades, usuários, níveis de acesso, assinatura, tramitação, blocos e histórico de movimentação, sem confundir ferramenta com o regime jurídico do processo administrativo.

\section{Síntese conceitual}

\mapafuncional{Competências Digitais}
{Competência digital = compreender informação + operar tecnologia + comunicar + agir com segurança e responsabilidade}
{Cultura digital $\rightarrow$ letramento e pensamento computacional.\\Internet/Web $\rightarrow$ busca, navegador, e-mail e colaboração.\\Segurança $\rightarrow$ autenticação, malware e backup.\\Direitos digitais $\rightarrow$ LGPD, LAI, Marco Civil e PNED.\\Setor público $\rightarrow$ interoperabilidade e processo eletrônico.}
{HTTPS protege canal, não a veracidade do site.\\Sincronização mantém estado; backup preserva capacidade de recuperação.\\MFA exige fatores distintos.\\LGPD: identifique dado, finalidade, base, agente e direito do titular.\\Interoperabilidade exige mais que conectividade.}
{Modo anônimo $\neq$ anonimato.\\Nuvem $\neq$ informação pública.\\Consentimento $\neq$ base universal.\\Cookie $\neq$ malware.\\Fonte mais ranqueada $\neq$ fonte mais confiável.}
{Qual é a fonte primária?\\Que dado está sendo tratado?\\Qual é a finalidade?\\Que identidade pode acessar?\\Há cópia recuperável?\\O sistema troca apenas bytes ou também significado?}

\begin{resumo}
Este capítulo deve ser estudado como integração entre informática básica, segurança, cidadania digital e regime jurídico do uso de informação. A cobrança mais seletiva costuma apresentar uma situação concreta e exigir que o candidato diferencie conceitos próximos: Internet e Web, sincronização e backup, autenticação e autorização, dado pessoal e dado anonimizado, transparência e exposição indevida.
\end{resumo}
